% -*- coding: utf-8 -*-
% !TEX root = ../main.tex

\subsection{Périmètre et sources de l'analyse}
\label{ssec:ddrs_perimetre}

Cette annexe porte sur l'ENSIIE, établissement d'accueil du stage, et plus
précisément sur le C19, où le projet a été réalisé. Elle s'appuie sur les
documents publics de l'école : la page consacrée à sa démarche DD\&RS, son
rapport d'activité DD\&RS 2022--2023, l'annonce du renouvellement de son label
en 2025 et ses indicateurs sociaux. Ces documents sont complétés par des
publications de l'ADEME, de l'Arcep, de l'Agence internationale de l'énergie et
des ministères compétents. Ils permettent de distinguer les engagements
institutionnels, les réalisations documentées et mon analyse personnelle de
leurs limites \cite{ensiieDDRS2026,ensiieRapportDDRS2023,ensiieLabelDDRS2025}.

\subsection{Pressions exercées sur l'environnement}
\label{ssec:ddrs_pressions}

\subsubsection{Enseignement supérieur, recherche et secteur numérique}

Une école d'ingénieurs en informatique exerce d'abord les pressions communes à
tout établissement d'enseignement supérieur : chauffage et électricité des
bâtiments, déplacements quotidiens et professionnels, restauration, achats et
déchets. Sa spécialisation ajoute une dépendance forte aux équipements
électroniques, aux réseaux et aux centres de données. Les impacts ne se
limitent donc pas à l'électricité consommée pendant un calcul. La fabrication
des ordinateurs et serveurs mobilise des métaux, des minerais, de l'eau, des
plastiques et de l'énergie ; leur extraction et leur transformation contribuent
aux émissions de gaz à effet de serre, à la pollution, à l'artificialisation
des sols et à la dégradation d'écosystèmes. La fin de vie crée enfin des déchets
électriques et électroniques dont le traitement est complexe.

L'étude ADEME--Arcep estime qu'en 2022 le numérique représentait en France
$29{,}5~\mathrm{MtCO_2e}$, soit $4{,}4\,\%$ de l'empreinte carbone nationale,
et $51{,}5~\mathrm{TWh}$, soit $11\,\%$ de la consommation électrique. Elle
attribue environ $50\,\%$ de l'empreinte carbone aux terminaux, $46\,\%$ aux
centres de données et $4\,\%$ aux réseaux. La production des équipements
mobiliserait en outre 117 millions de tonnes de ressources par an. Ces ordres
de grandeur montrent que la prolongation de la durée de vie du matériel est au
moins aussi importante que l'optimisation de son usage
\cite{ademeEmpreinteNumerique2025}.

L'intelligence artificielle accentue la pression sur les centres de données.
À l'échelle mondiale, ceux-ci ont consommé environ $415~\mathrm{TWh}$ en 2024,
soit $1{,}5\,\%$ de l'électricité mondiale. Le scénario central de l'Agence
internationale de l'énergie atteint environ $945~\mathrm{TWh}$ en 2030, l'IA
constituant le principal moteur de cette croissance avec les autres services
numériques \cite{aieEnergyAI2025}. Cette hausse comporte également des enjeux
locaux : puissance du réseau électrique, eau utilisée directement ou
indirectement pour le refroidissement, emprise foncière et conflits d'usage.

Les estimations de consommation par requête doivent être utilisées avec
précaution. Le rapport de l'Arcep publié en 2026 donne, selon les modèles et les
méthodes de calcul, des valeurs comprises entre environ $0{,}3$ et
$4~\mathrm{Wh}$ par requête. La longueur de la réponse, le matériel,
l'architecture et le lieu d'exécution peuvent changer fortement le résultat.
Le même rapport compare aussi l'entraînement de deux grands modèles de taille
proche. Le premier demande $1\,287~\mathrm{MWh}$ et produit
$552~\mathrm{tCO_2e}$, contre $433~\mathrm{MWh}$ et $30~\mathrm{tCO_2e}$ pour
le second. Cet écart montre que le nombre de paramètres ne suffit pas pour
mesurer l'impact environnemental \cite{arcepIAEnvironnement2026}.

\subsubsection{Spécificités de l'ENSIIE et du C19}

À l'ENSIIE, les postes de travail, les salles d'enseignement, les serveurs et
les déplacements constituent les principales sources directes ou indirectes
de pression attendues. Le C19 concentre certaines d'entre elles : l'école y
présente un centre de données, un parc informatique, des équipements de
réalité virtuelle et augmentée, deux salles immersives, ainsi que des espaces
de travail pour des projets et des jeunes entreprises. Ces ressources sont
cohérentes avec une activité de formation et d'innovation, mais leur
fabrication, leur renouvellement et leur utilisation doivent être inclus dans
l'empreinte du site \cite{ensiieC19}.

\subsubsection{Spécificités de l'équipe et de la mission}

La mission a mobilisé un poste CPU, du stockage local et une succession
d'entraînements et d'évaluations. Elle reste sans commune mesure avec
l'entraînement d'un grand modèle de langage : le corpus contient quelques
dizaines de milliers de textes, le backbone CBOW est de petite taille et les
expériences finales ont été exécutées sans GPU.

\subsection{Cadre réglementaire et institutionnel}
\label{ssec:ddrs_reglementation}

Plusieurs textes encadrent ces pressions, avec des niveaux d'obligation
différents.

Le dispositif Éco Énergie Tertiaire, issu de la loi Élan, concerne les
bâtiments ou ensembles de bâtiments accueillant au moins $1\,000~\mathrm{m^2}$
d'activités tertiaires, notamment l'enseignement et les salles informatiques.
Il fixe une réduction de la consommation d'énergie finale d'au moins
$40\,\%$ en 2030, $50\,\%$ en 2040 et $60\,\%$ en 2050 par rapport à 2010,
avec un suivi annuel sur la plateforme OPERAT
\cite{ecologieEcoEnergieTertiaire2025}.

Pour le numérique, la loi visant à réduire son empreinte environnementale a
conduit au Référentiel général d'écoconception des services numériques. Le
RGESN propose 78 critères portant notamment sur l'utilité du service,
l'architecture, l'hébergement, la limitation des ressources et la transparence
environnementale. Il constitue un cadre opérationnel pertinent pour les
logiciels et prototypes de recherche, même lorsque ceux-ci ne sont pas encore
des services déployés \cite{arcepRGESN2024}.

\subsection{Politique et actions de l'ENSIIE}
\label{ssec:ddrs_politique_ensiie}

\subsubsection{Gouvernance et communication}

L'ENSIIE s'est dotée d'une charte DD\&RS en 2019. Sa démarche est portée par
la direction, une personne chargée de mission et un comité de pilotage associant
plusieurs parties prenantes. Elle couvre les cinq axes du référentiel national :
stratégie et gouvernance, formation, recherche, environnement et politique
sociale. L'école a obtenu le renouvellement de son label DD\&RS en 2025 pour
quatre ans. L'audit met notamment en avant la gouvernance, l'implication des
élèves, la recherche liée à la transition, la végétalisation, la démarche bas
carbone et la politique sociale \cite{ensiieLabelDDRS2025}.

\subsubsection{Fonctionnement du campus}

La charte prévoit de rationaliser l'énergie, l'eau et les matériaux, de
poursuivre la rénovation énergétique, de réaliser des diagnostics écologiques,
d'adopter des achats durables, d'organiser le tri par filières et de travailler
sur les déplacements. Parmi les réalisations publiées figurent un bilan carbone
portant sur 2018, un plan de mobilité inter-établissements, des poubelles de tri,
un atelier de réparation informatique et la revégétalisation du square. Ces
actions couvrent à la fois la sobriété d'usage, la durée de vie du matériel et
la biodiversité locale \cite{ensiieDDRS2026,ensiieRapportDDRS2023}.

\subsubsection{Formation, recherche et sensibilisation}

L'intégration au cœur de métier est plus structurante que les seuls éco-gestes.
Le rapport 2022--2023 indique que tous les élèves ont suivi un enseignement
DD\&RS depuis 2019 et que 170 élèves ont participé à la rentrée climat. Il cite
des Fresques du climat, des ateliers 2tonnes, un cours sur les impacts
environnementaux du numérique et des travaux menés avec la ville
d'Évry-Courcouronnes. Du côté de la recherche, onze enseignants-chercheurs
avaient déclaré des thématiques liées au DD\&RS ; l'école participait également
à Labos 1point5 et au GDS EcoInfo. Le C19 relie enfin recherche, innovation et
territoire, ce que l'audit de 2025 identifie comme un point fort
\cite{ensiieRapportDDRS2023,ensiieLabelDDRS2025}.

Les documents établissent donc l'existence d'une offre de sensibilisation. En
revanche, aucune session DD\&RS propre à la courte période de ce stage
n'apparaît dans le planning ou les artefacts du projet. Pour cette mission, la
sensibilisation a surtout été indirecte, par l'étude du compromis entre qualité
et ressources. Ce décalage constitue un point d'amélioration : un module bref
d'accueil, accessible aussi aux stagiaires et intervenants temporaires,
permettrait de présenter la politique du site, les consignes locales et la
méthode de mesure attendue pour les projets numériques.

\subsection{Lien entre la mission et la stratégie environnementale}
\label{ssec:ddrs_lien_mission}

La compression étudiée relève de l'écoconception algorithmique : réduire les
paramètres peut limiter mémoire et calcul en inférence
\cite{arcepIAEnvironnement2026}. Les masques OBD et MINT suppriment environ la
moitié des poids de H2, mais le format dense conserve la même taille et la même
latence. À l'inverse, l'architecture compacte M4 réduit la taille du pipeline de
$34{,}95\,\%$ et sa latence bout en bout de $12{,}83\,\%$.

Le coût de calcul doit aussi être compté : le scoring H2 demande $1{,}41$~s avec
OBD contre $230{,}1$~s avec MINT. Comme aucune énergie n'a été mesurée, ces
résultats décrivent un potentiel de réduction, et non une économie en
kilowattheures ou en émissions.

\subsection{Responsabilité sociétale et usages du modèle}
\label{ssec:ddrs_responsabilite}

\subsubsection{Données, vie privée et biais}

Le corpus Hugging Face est sous licence Apache~2.0, mais sa provenance, son
annotation et sa population sont peu documentées
\cite{huggingfaceSentiment2023}.

Le RGPD reste applicable aux données publiques : les fichiers doivent être
contrôlés, minimisés et conservés pendant la seule durée utile
\cite{cnilDonneesInternet2024}.

La conversion en classification binaire et l'absence de métadonnées empêchent
un audit fiable des biais ; un déploiement réel nécessiterait des données
représentatives, une documentation des erreurs et une supervision humaine.

\subsubsection{Politique sociale de l'établissement}

La politique sociale de l'ENSIIE couvre égalité, diversité, accessibilité et
prévention des violences. Les indicateurs disponibles mentionnent $20\,\%$ de
femmes parmi les élèves, $18\,\%$ parmi les diplômés 2024--2025 et un écart de
salaire moyen d'environ $4{,}5\,\%$ en 2023
\cite{ensiieIndicateursDDRS2026}.

Ces indicateurs rendent le suivi possible, mais confirment que la féminisation
et l'inclusion restent des objectifs à poursuivre.

\subsection{Analyse critique et leviers d'amélioration}
\label{ssec:ddrs_critique}

L'ENSIIE dispose d'une gouvernance DD\&RS structurée et d'actions dans la
formation, la recherche et le campus.

La limite principale est l'absence de séries publiques récentes sur l'énergie,
l'eau, les mobilités, les achats et le numérique, qui empêche de mesurer les
effets réels des actions.

Trois leviers permettraient d'aller plus loin :

\begin{enumerate}
    \item \textbf{Mesurer et publier une trajectoire.} Un bilan annuel par
    grands postes, rapporté aux usagers et aux surfaces, rendrait les progrès
    vérifiables.

    \item \textbf{Piloter le numérique du C19.} La durée de vie, la réparabilité,
    l'énergie et la fin de vie du matériel et du centre de données devraient être
    suivies par des indicateurs.

    \item \textbf{Documenter les projets numériques.} Une fiche devrait regrouper
    ressources, données, essais, usages, biais et point d'amortissement, et
    vérifier les gains avec un format réellement compact.

\end{enumerate}

Le projet confirme qu'une sparsité théorique ne devient un gain environnemental
que si le matériel ou le moteur d'exécution l'exploite ; une mesure énergétique
directe reste nécessaire.
